\section{Discussion}
\label{sec:discussion}

% OWNER: Sanjana Garimella (results/discussion per proposal)

\paragraph{When does the rule break?}
Two findings constrain when subject--verb agreement survives noisy training.
First, the breakdown is \emph{threshold-like}, not gradual: accuracy is nearly
flat up to a $10\%$ corruption rate and only collapses at $50\%$
(Table~\ref{tab:main}). Small amounts of ungrammatical input are absorbed, much
as a child's grammar is not derailed by occasional speech errors. Second, the
attractor gap is a more sensitive early-warning signal than overall accuracy:
it climbs steadily ($0.073 \to 0.198$) even while aggregate accuracy looks
healthy, revealing that the model is quietly trading the hierarchical rule for
the linear nearest-noun heuristic well before its overall performance gives way.
The trajectory analysis sharpens the distinction between the two ways noise can
act. Up to $10\%$ noise, learning is merely \emph{delayed}---every condition
converges to the same accuracy given enough steps---whereas at $50\%$ it is
\emph{prevented}: the model plateaus far short of the rule, and further training
deepens rather than repairs its reliance on the linear cue.

\paragraph{Does structure buy robustness?}
Our central architectural question---whether a syntactically biased model (RNNG)
withstands noise better than a purely sequential one (LSTM)---the RNNG sweep now
answers in the affirmative (Section~\ref{sec:rnng}). At matched noise the RNNG is
more accurate than the LSTM throughout, and the advantage \emph{grows} with
noise, from under two points on clean data to $12.6$ points at $50\%$
corruption. The decisive contrast is in \emph{which rule} each model keeps: the
LSTM's attractor gap climbs steeply as training data degrades
($0.073 \to 0.198$), marking a drift toward the linear nearest-noun shortcut,
whereas the RNNG's gap stays low and almost flat ($0.020 \to 0.036$)---under the
heaviest noise it is still only half the LSTM's clean-data gap. An explicit
hierarchical prior thus preserves the structural generalization at noise levels
where a sequential learner has already defaulted to the linear one, implicating
architectural structure---not just input statistics---as a source of robustness.
The finding is qualified: the RNNG still loses overall accuracy at $50\%$ noise
($0.979 \to 0.834$), so structure blunts rather than abolishes the threshold
effect, and these RNNG values rest on a single (full-scale) seed pending the
multi-seed replication. Read against the poverty-of-the-stimulus framing of
Section~\ref{sec:motivation}, the result cuts toward the nativist side of the
ledger: a structural bias measurably improves what is learnable from the same
degraded input.

\paragraph{Cognitive-modeling angle.}
The pattern is suggestive for how structural knowledge might arise from
statistical learning. The LSTM does not begin agreement-aware; the hierarchical
rule is \emph{acquired} over training, and in the clean setting the attractor
gap shrinks as that rule consolidates. Noise does not flip a switch so much as
shift the balance between two competing generalizations---hierarchical and
linear---both of which the data can support. Heavy enough noise makes the linear
rule the better fit to the (corrupted) input, and the model follows the data.
That a distributional learner recovers the hierarchical rule at all from this
templatic corpus is itself evidence that the structure is more learnable from
input than a strong nativist account would predict; that it abandons the rule
under sufficient noise marks the limit of that learnability.

\paragraph{Threats to validity.}
Three caveats bound these claims; we treat them fully in
Section~\ref{sec:limitations}. Our noise is synthetic and uniform---a fixed
per-sentence flip probability---rather than the structured, frequency-dependent
errors of natural input. The corpus is templatic and vocabulary-controlled,
which aids control but limits external validity. Finally, the trajectory
analysis is single-seed; the converging-versus-plateauing pattern is clear, but
the precise crossover noise rate should be read as approximate until replicated
across seeds.
